\section{Conclusion}
\label{sec:conclusion}

This survey reviewed \Ncore{} studies that measure or intervene on the internal representations of
medical vision, image--text and multimodal foundation models, organised by method family, by the locus
at which the representation is read, and by whether the evidence is observational or interventional.
Table~\ref{tab:synthesis} states what the corpus reports, what it stops short of, and the experiment
that would settle each question. The five conclusions are not equally supported. Decodability of
findings, anatomy, demographic and acquisition attributes rests on the widest spread of independent
groups in the corpus, \Fdecgrpmin{} to \Fdecgrpmax{} distinct first authors over four to seven
modalities, but on records that almost never test selectivity. Site and scanner structure is recovered
in every record that tests it but one, yet a comparison model appears in a minority of those records
and one modality dominates. The meaning of discovered features rests on \Nsae{} dictionaries, of which
\Nsaeval{} report a semantic validation and \Nsaestab{} any reproducibility check, so the negative half
of that conclusion is firmer than the positive half. The comparison of medical with general
pre-training is uncontrolled by construction, so its conclusion is an absence of identification rather
than a finding. And the dependence of decodability on where the representation is read is a hypothesis
from three heterogeneous studies and one successive-locus probe. One qualification applies to all five:
\Pctpreprint\% of the corpus is preprints, and the appraisal is a description of reported controls
(Fig.~\ref{fig:findings}), not a certainty grade. Three consequences follow for anyone who studies or
deploys a medical foundation model.

\emph{Breadth is not control.} The checks these studies run are not interchangeable, so the survey
counts them in \Vnclasses{} classes, each stated with what it licenses and what it does not
(Table~\ref{tab:vcov}). A statement that lacks a class is narrower, not refuted. A probe establishes
that a readout recovers the target, and no observational control turns that into use: a control task
adds selectivity and a random encoder adds attribution to training, each a weaker statement than use. A
dictionary feature is a clinical concept only after clinician or structured-label validation, however it
was named, and a steering effect is direction-selective only against matched controls; a matched control
is reported in \Vspc{} of \Vspcn{} targeted interventions, and no medical study in this corpus meets the
criterion of \S\ref{sec:methods} for a concept-specific causal contribution.

\emph{What to measure.} Recovery is reported broadly and use is barely measured, and reported broadly is
not the same as established: \Fsensrefemptyn{} statements lose every record when the analysis is
restricted to peer-reviewed studies, only \Nminvalid{} of the \Ncore{} studies state an independent
split, evaluate on unused data and report uncertainty, and controls on probes are rare. A new audit
should not add another probe. It should pair one with a matched intervention or an erasure at the same
locus, and report both.

\emph{Where to measure.} The readout, whether pooling, depth or the direction chosen, can change what
appears recoverable, and the connector output and text encoder of multimodal models remain almost
unmeasured as loci of their own; a locus should be named before a model is praised. Whether medical
pre-training itself changes any of this is unidentified in the comparisons that exist, and settling it
takes paired pre-training on matched corpora, not another convenience comparison.
